\section{Datos y corpus}
\label{sec:datos}

El corpus proviene de MSD Manual Topics, un conjunto público con textos de manuales médicos en inglés, español, francés y alemán \cite{piperno2025}. Cada lengua tiene versiones para pacientes y para profesionales, y las muestras se etiquetan con temas jerárquicos. Los autores conservaron las 18 clases comunes a las cuatro lenguas con al menos 400 muestras cada una y fijaron un máximo de 512 tokens por texto. Antes del filtrado había 81.300 muestras de entrenamiento y 27.100 de prueba; después quedaron 52.647 en los conjuntos de entrenamiento, de las cuales 47.382 se destinaron al ajuste y 5.265 a validación, y 17.490 de prueba. La Tabla~\ref{tab:corpus} resume los conteos por lengua.

\begin{table}[H]
\centering
\caption{Muestras por lengua después del filtrado por clases.}
\label{tab:corpus}
\small
\setlength{\tabcolsep}{4pt}
\begin{tabular}{@{}lrrrr@{}}
\toprule
Lengua & Entrenamiento & Validación & Prueba & Prueba filtrada \\
\midrule
Inglés  & 11.860 & 1.325 & 4.295 & No aplica \\
Español & 12.171 & 1.351 & 4.519 & 1.842 \\
Francés & 11.748 & 1.298 & 4.308 & 1.947 \\
Alemán  & 11.603 & 1.291 & 4.368 & 2.317 \\
\bottomrule
\end{tabular}

\vspace{4pt}
{\footnotesize Nota: elaboración propia a partir de la Tabla~2 de \cite{piperno2025}. La columna de entrenamiento excluye la validación. La prueba filtrada aplica el umbral de BLEU de 0,3 y solo existe para las lenguas destino; el inglés conserva la prueba original.}
\end{table}

La destilación necesita pares entre el inglés y cada lengua destino. Los autores tradujeron el entrenamiento inglés con mBART-large-50 \cite{liu2020} y filtraron la prueba con BLEU: descartaron las muestras cuya similitud con el entrenamiento traducido superaba 0,3. El filtro conserva el 40,8\,\% de la prueba española, el 45,2\,\% de la francesa y el 53,0\,\% de la alemana, de modo que elimina entre el 47 y el 59\,\% de las muestras. El artículo llama \textit{unbiased} al conjunto resultante, pero el término se refiere a la ausencia de solapamiento léxico, no a equidad. El corpus son manuales educativos, no historias clínicas; las cuatro lenguas son europeas y el artículo no especifica la variedad del español.
